
\definecolor{lstbg}{RGB}{250,250,248}
\definecolor{lstframe}{RGB}{210,210,210}
\definecolor{lstkey}{RGB}{0,92,175}
\definecolor{lststr}{RGB}{163,21,21}
\definecolor{lstcom}{RGB}{0,120,0}
\definecolor{lsttag}{RGB}{136,18,128}
\definecolor{lstnum}{RGB}{130,130,130}

\lstdefinestyle{appendixbase}{
  basicstyle=\ttfamily\fontsize{6.6}{7.4}\selectfont,
  backgroundcolor=\color{lstbg},
  frame=single,
  rulecolor=\color{lstframe},
  framerule=0.5pt,
  framesep=3pt,
  xleftmargin=0.15em,
  xrightmargin=0.15em,
  aboveskip=0.5\baselineskip,
  belowskip=0.4\baselineskip,
  belowcaptionskip=0.6\baselineskip,
  columns=fullflexible,
  keepspaces=true,
  showstringspaces=false,
  showtabs=false,
  tabsize=2,
  upquote=false,
  breaklines=true,
  breakatwhitespace=false,
  prebreak=\mbox{\textcolor{lstnum}{\tiny$\hookleftarrow$}},
  postbreak=\mbox{\textcolor{lstnum}{\tiny$\hookrightarrow$}\space},
  captionpos=t,
  numbers=none
}

\lstdefinestyle{appendixtext}{
  style=appendixbase
}

\lstdefinelanguage{xmlcompact}{
  language=XML,
  morecomment=[s]{<!--}{-->},
  commentstyle=\color{lstcom},
  stringstyle=\color{lststr},
  identifierstyle=\color{black},
  keywordstyle=\color{lsttag}
}

\lstdefinestyle{appendixxml}{
  style=appendixbase,
  language=xmlcompact
}

\lstdefinelanguage{jsoncompact}{
  basicstyle=\ttfamily\fontsize{6.6}{7.4}\selectfont,
  showstringspaces=false,
  breaklines=true,
  breakatwhitespace=false,
  morestring=[b]",
  stringstyle=\color{lststr},
  literate=
   *{0}{{{\color{lstnum}0}}}{1}
    {1}{{{\color{lstnum}1}}}{1}
    {2}{{{\color{lstnum}2}}}{1}
    {3}{{{\color{lstnum}3}}}{1}
    {4}{{{\color{lstnum}4}}}{1}
    {5}{{{\color{lstnum}5}}}{1}
    {6}{{{\color{lstnum}6}}}{1}
    {7}{{{\color{lstnum}7}}}{1}
    {8}{{{\color{lstnum}8}}}{1}
    {9}{{{\color{lstnum}9}}}{1}
    {:}{{{\color{black}{:}}}}{1}
    {,}{{{\color{black}{,}}}}{1}
    {\{}{{{\color{black}{\{}}}}{1}
    {\}}{{{\color{black}{\}}}}}{1}
    {[}{{{\color{black}{[}}}}{1}
    {]}{{{\color{black}{]}}}}{1},
  moredelim=**[is][\color{lstkey}]{@}{@}
}

\lstdefinestyle{appendixjson}{
  style=appendixbase,
  language=jsoncompact
}

\begin{lstlisting}[style=appendixtext,caption={System prompt (compacted for appendix display).},label={lst:appendix-system-prompt}]
--- SYSTEM MESSAGE ---

You are an expert GPU performance engineer and compiler developer.
Your task is to analyze a GPU kernel and accurately estimate its launch configuration, FLOP operations, and DRAM accesses during its FIRST invocation.

You will be provided with the following information in XML tags:
- <program_name>: The name of the benchmark program.
- <kernel_mangled_name> and <kernel_demangled_name>: The target GPU kernel to analyze.
- <command_line_input_args>: The benchmark execution arguments (exe_args).
- <gpu_roofline_specs>: Hardware specifications for the target GPU architecture.
- <compile_commands>: The compilation commands used, which define macros and include paths.
- <source_code>: The source code files required for analysis.
- <sass>: Optional hardware SASS instructions. When present, treat SASS as the strongest low-level evidence for what instructions actually execute on the taken path, including relevant helper functions or inlined code reached by the target kernel.

Estimate the following metrics for the FIRST invocation of the target kernel:
Section 1:
- CUDA Grid Size (gridSz) as a 3-tuple of integers (gridSz_explanation, gridSz)
- CUDA Block Size (blockSz) as a 3-tuple of integers (blockSz_explanation, blockSz)
Section 2:
- Half-precision (FP16) FLOP count (fp16_flop_explanation, fp16_flop_count)
- Single-precision (FP32) FLOP count (fp32_flop_explanation, fp32_flop_count)
- Double-precision (FP64) FLOP count (fp64_flop_explanation, fp64_flop_count)
Section 3:
- DRAM bytes read (dram_bytes_read_explanation, dram_bytes_read_count)
- DRAM bytes written (dram_bytes_written_explanation, dram_bytes_written_count)

... additional task-specific instructions omitted for brevity ...

---

## Output Format
For each metric, provide a short two-sentence explanation of how you arrived at the final count. Include the reasoning behind the total work or traffic performed in the first kernel invocation, along with any important assumptions or simplifications.
Remember this evidence hierarchy: source and host-code reasoning establish the high-level execution model, and SASS is the strongest low-level evidence for the executed path when present. State uncertainty explicitly whenever dynamic behavior cannot be derived from source or SASS.
\end{lstlisting}


\begin{lstlisting}[style=appendixxml,caption={Example XML-style human prompt used for static analysis (compacted for appendix display).},label={lst:appendix-human-prompt}]
--- HUMAN MESSAGE ---

<program_name>
    adam-cuda
</program_name>
<kernel_mangled_name>
    _Z4adamIffEvPT_S1_S1_PKT0_fffffim10adamMode_tf
</kernel_mangled_name>
<kernel_demangled_name>
    void adam<float, float>(T1 *, T1 *, T1 *, const T2 *, float, float, float, float, float, int, unsigned long, adamMode_t, float)
</kernel_demangled_name>
<command_line_input_args>
    10000 200 100
</command_line_input_args>

<gpu_roofline_specs>
{
  "gpu_target": "H100 (SXM5)",
  "architecture_family": "Hopper",
  "architecture_chip": "GH100",
  "arch": "sm_90",
  "compute_capability": "9.0",
  "sm_count": 132,
  "memory_type": "HBM3",
  "global_memory_size_gb": 80,
  "memory_bandwidth_gb_per_s": 3360,
  "l2_cache_size_mb": 50,
  "per_sm": {
    "l1_cache_size_kb": 256,
    "cuda_cores": 128,
    "shared_memory_bytes": 233472,
    "shared_memory_human": "up to 228 KB"
  }
}
</gpu_roofline_specs>

<compile_commands>
<compile_command filename="HeCBench/src/adam-cuda/main.cu">
/opt/nvidia/hpc_sdk/Linux_x86_64/25.11/cuda/13.0/bin/nvcc ... -DGPU -O3 -DNDEBUG -std=c++17 "--generate-code=arch=compute_90,code=[compute_90,sm_90]" -x cu -rdc=true -c .../HeCBench/src/adam-cuda/main.cu -o CMakeFiles/adam-cuda.dir/main.cu.o
</compile_command>
</compile_commands>

<source_code>
<file name="HeCBench/src/adam-cuda/reference.h">
typedef enum {
  ADAM_MODE_0 = 0, // eps under square root
  ADAM_MODE_1 = 1  // eps outside square root
} adamMode_t;
... omitted ...
</file>

<file name="HeCBench/src/adam-cuda/main.cu">
template <typename T, typename G>
__global__ void adam(
  T* __restrict__ p, T* __restrict__ m, T* __restrict__ v,
  const G* __restrict__ g,
  const float b1, const float b2, const float eps,
  const float grad_scale, const float step_size,
  const int time_step, const size_t vector_size,
  adamMode_t mode, const float decay)
{
  const size_t i = blockIdx.x * blockDim.x + threadIdx.x;
  const size_t totThreads = gridDim.x * blockDim.x;
  for (size_t j = i; j < vector_size; j += totThreads) {
    for (int t = 1; t <= time_step; t++) {
      T scaled_grad = g[j] / grad_scale;
      m[j] = b1 * m[j] + (1.f - b1) * scaled_grad;
      v[j] = b2 * v[j] + (1.f - b2) * scaled_grad * scaled_grad;
      float m_corrected = m[j] / (1.f - powf(b1, t));
      float v_corrected = v[j] / (1.f - powf(b2, t));
      float denom = (mode == ADAM_MODE_0)
        ? sqrtf(v_corrected + eps)
        : sqrtf(v_corrected) + eps;
      float update = (m_corrected / denom) + (decay * p[j]);
      p[j] -= step_size * update;
    }
  }
}

int main(int argc, char* argv[]) {
  ...
  const int vector_size = atoi(argv[1]);
  const int time_step   = atoi(argv[2]);
  const int repeat      = atoi(argv[3]);
  ...
  for (int i = 0; i < repeat; i++) {
    adam<float, float><<<grids, blocks>>>(
      d_p, d_m, d_v, d_g,
      beta1, beta2, eps, grad_scale, step_size,
      time_step, vector_size, mode, decay);
  }
  ...
}
</file>
</source_code>

<sass>
<section name="_Z4adamIffEvPT_S1_S1_PKT0_fffffim10adamMode_tf">
/*0000*/ LDC R1, c[0x0][0x28];
/*0010*/ S2R R29, SR_CTAID.X;
...
/*00a0*/ HFMA2.MMA R28, -RZ, RZ, 0, 0;
/*00b0*/ BSSY B9, `(.L_x_26);
... omitted ...
</section>

<section name="__cuda_sm3x_div_rn_noftz_f32_slowpath">
/*0000*/ SHF.R.U32.HI R3, RZ, 0x17, R5;
/*0010*/ BSSY B0, `(.L_x_14);
...
/*0070*/ ISETP.GT.U32.AND P0, PT, R0, 0xfd, PT;
... omitted ...
</section>

<section name="__cuda_sm20_sqrt_rn_f32_slowpath">
... omitted ...
</section>
</sass>
\end{lstlisting}

\begin{lstlisting}[style=appendixjson,caption={Example expected LLM tool-call response in JSON output format (compacted for appendix display).},label={lst:appendix-json-output}]
{
  @"gridSz_explanation"@: "vector_size = 10000, threadsPerBlock = 256. Grid = ceil(10000/256) = 40. So gridSz = (40, 1, 1).",
  @"gridSz"@: [40, 1, 1],
  @"blockSz_explanation"@: "The code sets threadsPerBlock = 256, so blockSz = (256, 1, 1).",
  @"blockSz"@: [256, 1, 1],
  @"fp32_flop_explanation"@: "The kernel launches 40*256 = 10240 threads, but only 10000 threads do useful work (grid-stride loop with vector_size=10000, totThreads=10240, so each thread processes exactly 1 element). Each active thread runs an inner loop of time_step=200 iterations. 
... omitted for brevity ...
Total per element: ~114 + 199*54 = 114 + 10746 = 10860\nTotal FP32 \u2248 10000 * 10860 = 108,600,000",
  @"fp32_flop_count"@: 108600000,
  @"fp64_flop_explanation"@: "No double-precision floating point operations are present in this kernel. All arithmetic is single-precision (float).",
  @"fp64_flop_count"@: 0,
  @"fp16_flop_explanation"@: "The SASS contains HFMA2.MMA instructions used for constant materialization. In the first iteration path: HFMA2 at 0x00a0, 0x0510, 0x0e70 = 3 HFMA2 instructions (each = 4 FP16 FLOPs). In the loop iterations (t=2..200, 199 iterations): HFMA2 at 0x1560 and 0x1c10 = 2 HFMA2 per iteration. Total HFMA2 per element = 3 + 199*2 = 401. Each HFMA2 = 4 FP16 FLOPs. Total FP16 FLOPs per element = 401*4 = 1604. Total = 10000 * 1604 = 16,040,000.",
  @"fp16_flop_count"@: 16040000,
  @"dram_bytes_read_explanation"@: "Each active thread (10000 threads) reads 4 arrays: g[j] (4 bytes, read once via LDG.E.CONSTANT), m[j] (4 bytes), v[j] (4 bytes), p[j] (4 bytes). These are read once at the start of the time_step loop. Total unique read footprint = 4 * 10000 * 4 = 160,000 bytes. The L2 cache is 50 MB, so this easily fits. However, within the inner loop (200 iterations), m[j], v[j], and p[j] are read and written repeatedly but stay in L1/L2 cache. The g[j] array is read once. The initial compulsory misses bring all data from DRAM. Total DRAM reads \u2248 160,000 bytes (unique footprint of all 4 arrays, each 40,000 bytes).",
  @"dram_bytes_read_count"@: 160000,
  @"dram_bytes_written_explanation"@: "Each active thread writes to m[j], v[j], and p[j] after the inner loop completes (3 STG.E instructions visible in SASS at 0x1ed0, 0x1ee0, 0x1ef0). The unique write footprint = 3 * 10000 * 4 = 120,000 bytes. This is far smaller than the 50 MB L2 cache, so all writes will remain in L2 as dirty cache lines and will not be evicted to DRAM during this kernel invocation. Therefore DRAM bytes written \u2248 0.",
  @"dram_bytes_written_count"@: 0
}
\end{lstlisting}
